% !TEX program = xelatex
% 独立编译：xelatex outline.tex → bibtex outline → xelatex outline.tex（两次）。
% 组织形式参照 Survey_LLMS_Robots_Programming/outline.tex。
% 数值依据截至 2026-10-14 核对的 72 场景实验数据；物理结果源为 physics_validation_20261013/analysis/ 与 platform_comparison_20261014/analysis/。
% 基线文献见 references.bib；逐条核验记录紧邻对应条目。
\documentclass[UTF8,fontset=fandol,a4paper,11pt]{ctexart}
\usepackage[margin=2.2cm]{geometry}
\usepackage{xcolor,booktabs,tabularx,array,longtable}
\usepackage{amsmath,amssymb,enumitem,needspace,placeins,float}
\usepackage{tikz}
\usetikzlibrary{arrows.meta,positioning}
\usepackage{xurl}
\usepackage[colorlinks=true,linkcolor=blue!55!black,urlcolor=blue!55!black]{hyperref}
\usepackage[capitalize,nameinlink]{cleveref}
\newcommand{\qyb}[1]{\textcolor{red!75!black}{#1}}
\pdfstringdefDisableCommands{\def\qyb#1{#1}}
\crefname{section}{节}{节}
\Crefname{section}{节}{节}
\crefname{table}{表}{表}
\Crefname{table}{表}{表}
\crefname{figure}{图}{图}
\Crefname{figure}{图}{图}
\newcolumntype{Y}{>{\raggedright\arraybackslash}X}
\newcolumntype{P}[1]{>{\raggedright\arraybackslash}p{#1}}
\setlist{leftmargin=2em,itemsep=.2em,topsep=.3em,parsep=0pt}
\setlength{\parskip}{.3em}
\setlength{\emergencystretch}{2em}
\renewcommand{\arraystretch}{1.16}
\newcommand{\rqone}{不同双目感知方法生成的工件点云质量如何？}
\newcommand{\rqtwo}{固定 GraspGen 时，不同双目感知方法的抓取效率如何？}
\newcommand{\rqthree}{固定 FoundationStereo 时，不同抓取模型的抓取成功率如何？}
\title{基于双目三维感知的工业工件抓取实证研究\\[.3em]
\large 论文提纲与实验结果组织方案\\[.4em]
An Empirical Study of Stereo Perception and\\
Grasp Generation for Industrial Manipulation}
\author{yubin qu}
\date{2026年10月14日}

\begin{document}
\maketitle
\pagestyle{plain}

本提纲围绕抓取前的双目图像、三维点云、六自由度抓取位姿、运动路径及物理执行组织论文。
沿用英文章节标题和中文内容说明，每个研究问题按 Motivation、Analysis Method、Results and Analysis、Answer to RQ 展开。
论文采用模块替换的受控仿真实验，实验数据集包含 8 件工件、72 个场景；RQ2 和 RQ3 均在三个仿真实验环境、三种场景条件下评价。
三个环境为 FANUC＋Isaac Sim、Franka＋Isaac Sim 和 Franka＋MuJoCo；三种场景条件为单件、邻件和料框。
本文所称工业工件为程序化生成的工业形状代理。
评价重点是工件能否在规定执行条件下被夹起并保持；新增工件与真实硬件验证作为未来工作。

\begin{abstract}
摘要拟按应用问题、比较设计、主要结果和适用范围组织。
首先说明工业抓取依赖深度感知、抓取生成和运动执行的衔接，单独评价深度误差或候选数量难以定位整条流程的失败来源。
随后概括统一双目输入、点云接口、抓取工具与执行协议，比较七种双目方法和五种抓取模型组成的十一条配置。
实验结果覆盖 72 个场景、504 份点云和 792 个候选生成任务；三个环境在两组材料条件下共完成 1244 次完整机械臂物理执行。
可用于摘要的实例是：全部点云满足接口点数要求，但近表面点占比和可见表面覆盖存在差异；
FoundationStereo 与 GraspGen 组合在默认材料条件下的抓取成功数依次为 FANUC＋Isaac Sim 的 42/72、Franka＋Isaac Sim 的 52/72 和 Franka＋MuJoCo 的 17/72。
\end{abstract}

\section{\qyb{Introduction}}
\label{sec:intro}

基于摄像头的机械臂抓取通过观测工件的位置与几何形状，为零件上料、分拣和料框取件生成可执行的抓取动作。
在采用双目感知的系统中，标定后的左右图像用于估计深度并生成三维点云，抓取模型据此预测夹爪位置、姿态和开口，
运动规划与机械臂控制随后完成接近、夹紧和抬升。
相关研究分别推进了三维感知与抓取生成：SGM 通过半全局匹配估计视差\cite{hirschmuller2008}，
FoundationStereo 使用预训练模型进行跨场景双目深度估计\cite{wen2025}；
Contact-GraspNet 从观测点云中的接触点预测六自由度抓取\cite{sundermeyer2021}，
GraspGen 则采用扩散模型生成候选，并通过判别器进行评分\cite{murali2026}。
这些方法为从相机图像到机械臂抓取的模块化系统提供了可替换的算法组件。

面向工业工位的算法选型，仍需要系统的实证研究说明点云生成方法与抓取模型组合后的执行表现。
上述代表性工作分别提供了深度估计或抓取性能证据\cite{wen2025,sundermeyer2021,murali2026}，
但仅凭各自在不同对象、传感器和机器人配置上的评价，尚不足以回答同一工业工位中替换某个模块后，最终抓取成功率会怎样变化。
工件的孔、棱边、邻件遮挡和料框边界会同时改变可见几何与机械臂接近空间，
点云误差、目标提取、抓取候选和运动约束也会共同作用于执行结果。
因此，本文关注的研究空白是统一工业工件、相机输入和执行规则下的模块比较证据：
不同方法生成的点云质量如何，这些差异怎样对应下游抓取结果，以及固定相同点云时不同抓取模型的表现如何。

针对这一研究空白，本文开展面向工业工件的模块替换受控仿真实验。
实验使用轴套／隔套和安装块两类共 8 件工件，构建包含单件、邻件和料框条件的 72 个场景，
比较七种双目感知方法和五种抓取模型组成的十一条配置。
首先评价点云的表面误差、可见覆盖和目标来源比例；随后固定 GraspGen 比较双目方法，并固定 FoundationStereo 比较抓取模型。
抓取任务分别在 FANUC＋Isaac Sim、Franka＋Isaac Sim 和 Franka＋MuJoCo 中评价，
各环境内采用共同的工具、规划和控制设置，两个 Franka 环境进一步复用相同关节路径。
实验将每个场景的点云、候选、路径和物理执行结果关联起来，以全场景抓取成功率及阶段记录分析模块替换与环境配置下的结果差异。

本文的贡献包括工业工件数据、受控比较结果和可复核的实验材料。
数据部分提供 72 个场景的双目 RGB、相机标定、工件几何与离线评价真值，为相同输入下的方法比较提供依据。
比较部分连接点云质量与三个仿真实验环境中的抓取结果，说明点数充足的输入仍可能存在几何偏差，
并展示同一感知与抓取组合在不同环境配置中的成功率差异。
复现材料包含统一模块接口、模型与参数配置、逐任务输入输出关联、成功与失败轨迹及分析代码，
用于核对各阶段结果并定位需要进一步验证的失败原因。
代码与实验日志通过 GitHub 仓库共享\footnote{\url{https://github.com/warmjademe/Robotic-Grasping-BasedOn-TwoCamerasSensing}}。

本文其余部分组织如下：\cref{sec:background}介绍双目感知、抓取生成及相关评价工作；
\cref{sec:design}给出研究问题、实验数据、基线方法、执行环境与评价指标；
\cref{sec:rq1}至\cref{sec:rq3}依次报告点云质量、固定抓取模型的感知比较和固定感知方法的抓取模型比较；
\cref{sec:discussion}讨论对工业界和学术界的启示及有效性威胁；\cref{sec:conclusion}总结研究并说明未来工作。

\section{\qyb{相关工作}}
\label{sec:background}

\subsection{\qyb{Stereo Perception and Point Cloud Construction}}

本小节拟解释左右图像如何经对应点匹配、视差到深度的转换和坐标变换形成米制点云。
以同一轴套的可见轮廓说明深度误差、表面缺失和非目标点混入分别怎样改变夹爪接触区域。
传统方法以 StereoBM 和 SGM 的现有实现为对照，学习方法包括 PSMNet、GwcNet、RAFT-Stereo、IGEV-Stereo 和 FoundationStereo。
技术背景按匹配代价与聚合、代价体学习、迭代估计和预训练模型组织，
重点说明输入、输出、标定依赖与运行配置；具体机制和训练数据以原论文及实际加载权重为依据。

\subsection{\qyb{Grasp Poses, Motion Plans, and Physical Execution}}

本小节拟明确抓取位姿、机械臂路径和控制动作的区别。
六自由度抓取位姿表示夹爪在三维空间的位置与朝向，候选还包含开口和方法内部评分；
运动规划将该位姿转换为满足关节、接近方向和碰撞约束的运动；控制与接触决定工件是否随夹爪抬起。
机械臂 API 执行的是给定动作，抓取生成模块负责选择夹持位置与姿态，两者在系统中承担不同职责。
用同一工件的两种候选说明：候选都可能靠近表面，但只有其中一种满足夹爪开口与接近空间要求。

\subsection{\qyb{Grasp Generation and Empirical Evaluation}}

本小节拟比较 GPD、6-DOF GraspNet、Contact-GraspNet、GraspNet-Baseline 和 GraspGen 的输入要求、候选表达与评价设置。
逐篇核对对象点云或场景点云、相机信息、候选采样、评分、夹爪几何及后处理，说明统一适配器保留或改变了哪些条件。
GPD 的分类器与其他模型均采用预训练参数，并记录实际版本。
ASGrasp 的额外主动感知输入作为相关工作的比较内容。
抓取数据集与仿真评价研究将按输入真实性、工件分布、机械臂约束、接触物理和成功定义比较，
据此确定本文与已有实验的对应关系。

\subsection{\qyb{研究空白}}
\label{sec:research-gap}

已有研究为三维感知和机器人抓取提供了方法与评价依据，但各项评价覆盖的系统环节不同。
FoundationStereo 比较双目深度估计，并展示真实场景中的点云结果，包括螺母和螺栓等对象\cite{wen2025}；
GraspNet-1Billion 提供真实 RGB-D 场景、抓取标注和统一的抓取位姿评价\cite{fang2020}；
GraspGen 则在仿真和真实机器人上验证抓取生成方法\cite{murali2026}。
这些工作已研究感知或抓取的实际表现，本文进一步关注这些模块组合后在工业机械臂任务中的执行结果。

工业工位的方法选择需要知道：同一批工件和摆放条件下，更换感知或抓取模块后，机械臂最终能否完成抓取。
从上述代表性工作的评价范围看，仍需补充统一工业工件、相机输入、夹爪和执行规则下的受控比较，
将点云误差与覆盖、抓取候选、路径可执行性和物理抓取成功逐任务关联起来。
单件、邻件和料框条件下的可见表面与接近空间不同，方法的离线几何或候选评价也需要与这些条件下的执行结果对应。
这类证据有助于判断失败发生在感知、抓取生成、运动规划还是夹持阶段，并支持工业系统的模块选择与调整。

本文通过点云质量评价及两组模块替换实验补充上述比较证据。
两组实验分别固定 GraspGen 比较感知方法、固定 FoundationStereo 比较抓取模型，
在三个仿真实验环境中覆盖单件、邻件和料框条件，以规定任务的抓取成功率评价整条执行流程。
当前实验使用合成工业形状和物理仿真，先检查模块组合在受控条件下的表现；
真实工业工位中的金属表面、相机噪声、机械臂标定及夹持参数将在后续硬件实验中验证。

\section{\qyb{Empirical Study Design}}
\label{sec:design}

本节拟使实验配置能够独立复现，集中说明任务、数据、模块、评价器和统计单位。

\subsection{\qyb{Study Scope}}

机器人操控可通过显式三维几何、世界模型或视觉--语言--动作（vision-language-action，VLA）模型组织感知与动作决策。
显式几何路线由相机获取深度或估计双目视差，构建点云，再生成抓取位姿并调用运动规划与控制模块\cite{wen2025,mousavian2019}；
世界模型学习动作引起的环境变化，例如 DayDreamer 利用预测轨迹学习控制策略，并在真实机械臂上完成视觉拾放\cite{wu2023daydreamer}；
VLA 模型则结合图像、语言指令和机器人状态生成动作，例如 Physical Intelligence 团队的 $\pi_0$ 通过流匹配输出连续动作\cite{black2025pi0}。
这些路线均可使用摄像头，其区别在于显式几何、环境变化预测与动作生成在系统中的作用。
本文选择基于相机深度与点云的抓取路线：Dex-Net 2.0 已在 ABB YuMi 上进行千余次实体抓取对照，GPD 也报告了杂乱场景中的真实机器人抓取实验\cite{mahler2017dexnet,tenpas2017}，
为这一路线的技术可行性提供了实验依据；同时，显式点云、抓取位姿和运动路径便于结合工业工位的标定、夹爪及避碰约束逐阶段检查。

本文研究工业工位中基于抓取前双目图像的静态工件抓取，流程如\cref{fig:pipeline}所示。
每次任务指定一个目标，由固定基座机械臂和平行两指夹爪完成接近、夹紧、抬升和保持。
例如，从相邻工件中拾取一个轴套时，左右相机先采集图像，感知方法估计深度并构建目标点云，
抓取模型据此生成夹爪位姿与开口，路径程序检查关节、接近和碰撞约束，再由机械臂执行选定动作。
实验覆盖单件、邻件和料框三种条件，分别观察工件几何、邻件遮挡与壁面约束下的抓取表现，
并将点云质量、候选、路径和物理执行记录关联起来，分析替换感知或抓取模块后的结果变化。

\begin{figure}[H]
\centering
\begin{tikzpicture}[
box/.style={draw,rounded corners=2pt,fill=gray!5,text width=4.35cm,minimum height=1.0cm,align=center,font=\small,inner sep=5pt},
note/.style={font=\small,align=center},>=Stealth]
\node[box] (rgb) at (0,0) {抓取前左右 RGB\\固定相机标定};
\node[box] (points) at (5.35,0) {深度与三维点云\\几何质量与目标提取};
\node[box] (grasp) at (10.7,0) {候选夹爪位姿与开口\\模型输出、坐标适配};
\node[box] (plan) at (10.7,-2.0) {共同路径检查\\关节、接近与碰撞约束};
\node[box] (physics) at (5.35,-2.0) {完整机械臂物理执行\\接近、闭合、抬升与保持};
\node[box] (results) at (0,-2.0) {逐任务阶段记录\\失败阶段与抓取成功率};
\draw[->] (rgb) -- (points);
\draw[->] (points) -- (grasp);
\draw[->] (grasp) -- (plan);
\draw[->] (plan) -- (physics);
\draw[->] (physics) -- (results);
\node[note] at (5.35,-3.25) {分别在三个环境和三种场景条件下比较\\FANUC＋Isaac Sim、Franka＋Isaac Sim、Franka＋MuJoCo};
\end{tikzpicture}
\caption{实验流程及评价职责。抓取模型给出位姿，路径程序给出运动；完整物理执行独立记录接触、抬升与禁止碰撞。}
\label{fig:pipeline}
\end{figure}

\subsection{\qyb{Research Questions}}

本文从几何输入质量和两类模块的替换效果出发，提出三个研究问题。

\begin{itemize}[leftmargin=3.2em,itemsep=.6em]
\item[RQ1.] \rqone

点云决定抓取模型能够利用的工件几何信息，因此首先比较各双目方法的点云质量。
固定左右图像、相机标定和目标提取步骤，从点数、表面误差、可见覆盖和目标来源比例分析差异，
为解释后续抓取表现提供逐场景的感知证据。
\item[RQ2.] \rqtwo

该问题考察感知方法的替换如何对应最终抓取表现，“抓取效率”以完成既定任务的成功比例衡量。
固定 GraspGen 及后续规划与控制设置，仅替换双目感知方法，比较全场景及不同摆放条件下的抓取成功率，
并结合候选生成、路径检查和物理执行记录，分析点云差异与下游失败阶段的关系。
\item[RQ3.] \rqthree

该问题考察相同几何输入下不同抓取模型的实际表现。
固定 FoundationStereo 输出及共同的夹爪、规划和控制设置，比较各抓取模型的成功率，
结合候选分布、适配和筛选记录，分析不同抓取生成方法在工业工件上的可用性。
\end{itemize}

\subsection{\qyb{实验数据集}}

数据集面向工业机械臂在工位上进行零件上料、分拣和料框取件的应用场景，评价平行两指夹爪对刚性工件的抓取（见\cref{tab:industrial-scenes}）。
实验数据集包含轴套／隔套和安装块两类程序化工业形状，每类四个尺寸变体，共 8 件工件。
工件包含孔、平面和棱边等工业零件常见结构，用于观察几何误差、邻件遮挡及料框边界对感知和执行的影响。
每件设置三个位姿，每个位姿包含单件、邻件和料框三种条件，共 $8\times3\times3=72$ 个场景，每种条件各 24 个。
每场景输入一对左右 RGB 图像，共 144 张；七种感知方法和三个仿真实验环境使用相同的 72 个场景清单。
渲染输入为 $640\times480$ RGB，双目基线为 60 mm、水平和垂直焦距均为 700 像素，
采用无畸变的合成针孔相机，感知输入为固定左右图像及其标定参数。
真值网格、对象身份、深度和位姿用于离线评价及物理环境构造。

\begin{table}[H]
\centering\small
\caption{面向工业工位的工件抓取评测数据集。每类场景均包含轴套／隔套和安装块两类共 8 件工件、每件 3 个位姿。工件为程序化工业形状，场景采用合成双目图像和物理仿真。}
\label{tab:industrial-scenes}
\begin{tabularx}{\linewidth}{P{1.8cm}P{2.4cm}Y r}
\toprule
场景类型 & 对应工业作业 & 工件摆放与抓取约束 & 场景数 \\
\midrule
单件 & 台面工件上料 & 目标工件单独放置于台面，比较不同位姿下的表面恢复、夹持位置选择及台面接近约束。 & 24 \\
邻件 & 相邻工件分拣 & 目标旁存在其他工件，评价非目标点混入、邻件遮挡与夹爪接近时的避碰要求。 & 24 \\
料框 & 料框内工件取件 & 目标位于带壁面约束的料框场景，评价壁面对可见表面、夹爪姿态及接近空间的影响。 & 24 \\
\midrule
合计 & 三类工业作业条件 & 相同 8 件工件在三种条件下重复摆放；共 72 对双目图像，144 张 RGB。 & 72 \\
\bottomrule
\end{tabularx}
\end{table}

\subsection{\qyb{基线方法}}

基线分为点云基线与抓取基线，分别替换双目深度估计模块和六自由度抓取生成模块。
选择策略兼顾技术路线覆盖、接口适配性和复现条件：点云方法覆盖传统几何匹配到预训练神经网络，
抓取方法覆盖几何采样与评分、直接预测和生成模型；各方法能够接入双目图像或米制点云接口，
并具有公开实现及可获取的必要预训练参数，以便记录和复核实际配置。
点云基线读取同一对校正图像，以相同标定将深度反投影为三维点；抓取基线读取共同点云接口并输出夹爪候选位姿。
学习方法采用已归档的预训练权重进行推理。

感知比较形成七种双目方法与 GraspGen 的组合，抓取比较形成 FoundationStereo 与五种抓取模型的组合。
两组共享 FoundationStereo 与 GraspGen，因此共有 $7+5-1=11$ 条不同配置。
各配置围绕共同参照替换单个模块，用于比较固定另一模块时的替换效果。
共享点云、候选和物理执行通过唯一任务标识关联。

\subsubsection{点云基线}

点云基线覆盖局部匹配、半全局优化、代价体网络、迭代更新和使用大规模预训练的双目方法。
各方法均使用固定的点云反投影与目标提取步骤，因此比较对象是深度估计模块在共同接口下产生的工件点云。

\begin{enumerate}
\item StereoBM 使用局部窗口的图像匹配代价计算视差，不需要预训练权重。
本文采用 OpenCV 实现，以 Konolige 的区域相关双目方法为技术背景\cite{konolige1998}，
具体接口以 OpenCV 的 StereoBM 文档为准\footnote{\url{https://docs.opencv.org/4.x/d9/dba/classcv_1_1StereoBM.html}}。
选择它作为局部匹配的简单参照，观察基本图像对应关系能够支持怎样的点云与抓取结果。
\item SGM 沿多个方向聚合匹配代价，并对视差变化施加平滑约束\cite{hirschmuller2008}。
本文使用 OpenCV StereoSGBM 实现并记录其参数，选择它代表无需预训练权重的半全局匹配，与 StereoBM 的局部匹配形成对照。
\item PSMNet 通过空间金字塔池化提取多尺度上下文，再用三维卷积处理代价体\cite{chang2018}。
选择它代表利用三维卷积学习匹配代价的网络，作为传统方法与后续神经网络的比较参照。
\item GwcNet 将左右特征分组计算相关性，构造供后续网络处理的匹配代价体\cite{guo2019}。
选择它补充代价体的另一种构建方式，与 PSMNet 比较不同匹配表示下的工件表面恢复表现。
\item RAFT-Stereo 利用相关性信息和多层循环更新模块，逐步修正视差估计\cite{lipson2021}。
选择它代表迭代视差估计方法，与代价体卷积网络比较点云质量及后续抓取表现。
\item IGEV-Stereo 将几何与上下文信息编码到体表示中，为视差初始化和迭代更新提供依据\cite{xu2023}。
选择它代表几何编码与循环更新相结合的方法，与 RAFT-Stereo 形成迭代估计路线中的对照。
\item FoundationStereo 结合大规模合成双目训练、单目视觉先验和上下文建模，面向跨场景的零样本双目匹配\cite{wen2025}。
选择它评价大规模预训练模型在工业工件上的直接迁移表现，并将其输出作为比较抓取模型时的固定感知输入。
\end{enumerate}

\subsubsection{抓取基线}

抓取基线覆盖几何候选与学习评分、变分生成、接触点预测、直接位姿预测和扩散生成五类方案。
这些方法输出候选夹法；运动规划与机械臂控制由共同执行程序完成。

\begin{enumerate}
\item GPD 在点云局部几何上采样平行夹爪候选，再利用学习得到的分类器评分\cite{tenpas2017}。
当前模块使用作者的几何描述子与预训练分类参数。
选择它代表几何候选采样与学习评分相结合的方案，为后续位姿预测和生成模型提供参照。
\item 6-DOF GraspNet 利用变分自编码器生成候选，并通过评价网络评分和修正候选位姿\cite{mousavian2019}。
选择它代表基于潜变量采样的抓取生成方法，与 GraspGen 的扩散生成路线比较候选可用性和执行结果。
\item Contact-GraspNet 从观测点云中的接触点预测抓取方向、开口等参数，生成六自由度夹爪位姿\cite{sundermeyer2021}。
选择它代表以可见表面接触点为依据的抓取预测，考察局部观测几何如何对应可执行的夹持方案。
\item GraspNet-Baseline 使用 GraspNet-1Billion 论文及官方基线中的点云抓取预测网络，分解预测接近方向与抓取参数\cite{fang2020}。
选择它作为公开抓取基准中的直接位姿预测参照，评价该预测方式迁移到工业工件后的表现。
\item GraspGen 利用扩散模型生成抓取候选，并由经过 on-generator training 的判别器评分\cite{murali2026}。
选择它代表扩散生成与学习评分相结合的方案，复用公开权重，并将其作为比较双目感知方法时的固定抓取模块。
\end{enumerate}

各基线保留原始候选、分数和适配后的夹爪坐标；评分仅用于模型内部排序。
配置表记录仓库提交、权重哈希、输入分辨率、采样设置、夹爪坐标映射和推理设备。

\subsection{\qyb{实验环境}}

共同感知接口输出米制场景点云和目标区域点云，并保留相机到工位坐标变换。
目标提取使用固定三维区域、预测高度超过桌面 3 mm 及近邻离群过滤，每条抓取候选提供位姿、开口和评分。
共同路径程序限制夹爪朝下且偏离竖直不超过 $45^\circ$，开口为 2--80 mm；
依据预测点云检查两指接触区域，以及从初始位姿到预抓取点的关节插值、120 mm 直线接近和 150 mm 直线抬升。
环境点云采用 3 mm 占据体素，关节与直线采样间隔分别约为 $1^\circ$ 和 5 mm。
该程序没有绕障搜索，结果分别记录候选筛除、逆运动学失败、碰撞和预算超时。

三个仿真实验环境由机械臂与仿真器的组合定义，包含两种机械臂模型和两种物理仿真器，见\cref{tab:environments}。
Isaac Sim 环境采用 5.1 版本、Isaac Lab 2.3.2.post1 与 PhysX，分别加载 FANUC M-20iD/25 模型和原生 Franka Panda。
MuJoCo 环境采用 3.3.7 版本及 Menagerie 的 Franka Panda 模型；所有环境使用平行两指工具。
工件初始化后的位姿由重力、夹指和场景接触产生的动力学过程更新。
完整流程包括接近、闭合、抬升、保持、降回、松开和退出。

三个环境复用已归档的双目输出与抓取候选，FANUC 和 Franka 分别生成符合各自运动学约束的计划。
两个 Franka 环境执行相同关节路径；FANUC 与 Franka 的比较按场景、感知方法和抓取模型匹配，保留各自已选候选和路径。
各环境分别通过夹持正例、夹空、夹爪不闭合和障碍阻挡校验；1244 次执行均已有结果及轨迹。
机器人结构、原生碰撞资产、惯性、驱动和接触求解实现的差别作为完整环境配置记录。

\begin{table}[H]
\centering\small
\caption{三个仿真实验环境。每个环境均覆盖相同的 72 个场景。执行次数包含默认与较低摩擦条件，排除孤立夹爪测试及校准。}
\label{tab:environments}
\begin{tabularx}{\linewidth}{Y Y r r}
\toprule
环境 & 机械臂与工具 & 已选计划 & 物理执行次数 \\
\midrule
FANUC＋Isaac Sim & M-20iD/25 模型＋Panda 两指工具 & 206 & 412 \\
Franka＋Isaac Sim & Isaac 原生 Franka Panda & 208 & 416 \\
Franka＋MuJoCo & Menagerie Franka Panda & 208 & 416 \\
\bottomrule
\end{tabularx}
\end{table}

三个环境中每件物体质量统一为 0.1 kg，默认滑动摩擦系数设为 1.0，较低摩擦条件为 0.5；实际材料与接触参数随环境配置归档。
夹指驱动力上限为每指 20 N，指关节 PD 参数为 1000/100，时间步为 1/250 s。
FANUC 的臂关节 PD 为 $10^6/10^5$，两个 Franka 环境为 400/80；各环境沿用其归档的驱动和动作阶段设置。
Franka 前四个臂关节的力矩上限为 87 N\,m，后三个为 12 N\,m；Isaac 中的 Franka 原生高增益配置关闭机器人重力，MuJoCo 对 Franka 启用重力补偿，工件均受重力作用。
碰撞门禁每物理步计算，位姿与接触轨迹按 50 Hz 保存；环境版本、材料读回和模型来源纳入复现配置。

完整执行通过要求工件底部距原支承面至少 100 mm 并保持 2 s，接近末端误差不超过 2 mm 和 $0.5^\circ$，
并且除初始化稳定阶段外的完整动作流程没有超过 0.5 N 的禁止接触合力。
物理抬升保持、接近到位和禁止接触分别保留布尔结果，联合判定作为执行通过。
放回精度另列：相对稳定初始位置的平面误差不超过 20 mm，工件底部距台面不超过 3 mm。
当前禁止接触使用非夹指连杆合力及夹指非目标接触合力，逐接触点覆盖的限制见\cref{sec:threats}。

历史 GraspDataGen 评价保持原候选相对位姿和已保存的闭合指令，在无桌面、关闭重力的条件下施加五段 1 g 拉拽。
该测试报告抗扰动夹持通过率，完整机械臂测试报告任务执行结果。
孤立夹爪测试只显式改变对象材料，夹爪沿用作者 USD 材料；完整机械臂材料条件覆盖机器人和场景的已读回碰撞形状。
两类测试按各自的外力、支承和材料条件并行报告诊断结果。

\subsection{\qyb{评价指标}}

评价指标分别描述几何输入质量和抓取任务完成情况。
以下以 $i$ 表示场景，以 $\mathbf{1}[\cdot]$ 表示条件成立时取 1、否则取 0 的指示函数；
距离以毫米报告，比例以百分数报告。

\subsubsection{点云质量}

深度误差衡量目标像素上的预测深度与真实深度之差。
设 $\Omega_i$ 为左图中具有有效真值深度的目标像素集合，$V_i$ 为方法输出的有效深度像素集合，
$D_i=\Omega_i\cap V_i$；$\hat z_i(u)$ 和 $z_i^*(u)$ 分别为像素 $u$ 的预测深度与真值深度。
平均绝对深度误差和有效深度比例定义为
\[
 E_i^{\mathrm{depth}}=\frac{1}{|D_i|}\sum_{u\in D_i}|\hat z_i(u)-z_i^*(u)|,
 \qquad
 v_i^{\mathrm{depth}}=\frac{|D_i|}{|\Omega_i|}.
\]
误差仅在有预测值的目标像素上计算，因此必须同时报告有效深度比例；$D_i$ 为空时误差记为缺测。
深度误差的 P95 为上述绝对误差样本的第 95 百分位数。

三维表面误差衡量输入抓取模型的目标点云偏离真实工件表面的程度。
设 $\mathcal P_i=(\mathbf p_{ij})_{j=1}^{n_i}$ 为经过接口校验的目标点序列，$n_i$ 为点数，$\mathcal S_i$ 为同一工位坐标系中的真实目标网格表面。
点到表面的最近距离及平均表面误差定义为
\[
 d(\mathbf p,\mathcal S_i)=\min_{\mathbf x\in\mathcal S_i}\|\mathbf p-\mathbf x\|_2,
 \qquad
 E_i^{\mathrm{surface}}=\frac{1}{n_i}\sum_{j=1}^{n_i}d(\mathbf p_{ij},\mathcal S_i).
\]
表面误差的 P95 为这组点到表面距离的第 95 百分位数；点数、非有限坐标数和重复点数分别记录。

近表面点占比和可见表面覆盖率分别衡量预测点的准确程度与观测表面的完整程度。
设 $\tau$ 为距离阈值，$\mathcal Q_i=(\mathbf q_{ik})_{k=1}^{m_i}$ 为左右相机真值目标深度反投影后拼接的表面采样点序列，则
\[
 p_i(\tau)=\frac{1}{n_i}\sum_{j=1}^{n_i}\mathbf{1}\!\left[d(\mathbf p_{ij},\mathcal S_i)\leq\tau\right],
 \qquad
 c_i(\tau)=\frac{1}{m_i}\sum_{k=1}^{m_i}\mathbf{1}\!\left[\min_j\|\mathbf q_{ik}-\mathbf p_{ij}\|_2\leq\tau\right].
\]
采样点按两个视角的有效目标像素拼接，不做面积加权或重叠表面去重，因此覆盖率针对左右可见表面的采样点。
阈值采用 1、2、5 mm，主表使用 2 mm。
目标点云为空时点数和两个比例记为 0，表面误差记为缺测；真值可见采样为空时覆盖率记为缺测。

目标来源比例衡量提取的目标点云中有多少点实际来自目标像素。
令 $M_i(u)$ 为左图真值目标掩码，$\pi_i(\mathbf p)$ 为使用标定参数将三维点投影并取整到左图的像素坐标，则
\[
 t_i=\frac{1}{n_i}\sum_{j=1}^{n_i}M_i\!\left(\pi_i(\mathbf p_{ij})\right).
\]
该比例与点到表面距离分别计算；非目标像素产生的预测点也可能落在工件表面附近。
空目标点云的来源比例记为缺测。真值掩码、深度和网格仅供离线评价使用。

\subsubsection{抓取成功率与阶段通过率}

抓取成功要求工件被抬起并保持、夹爪接近到位，且执行流程满足禁止接触约束。
设 $h_i(t)$ 为工件在时刻 $t$ 的最低网格顶点距原支承面的高度，$H_i$ 为规定保持阶段的采样时刻集合，
抬升保持标签为
\[
 L_i=\mathbf{1}\!\left[\min_{t\in H_i}h_i(t)\geq h_{\min}\right],
 \qquad h_{\min}=100\ \mathrm{mm}.
\]
保持阶段为 2 s，位姿记录频率为 50 Hz；该标签以完整保持记录为计算前提。
令 $\mathbf x_i,\mathbf R_i$ 为接近结束时的实际夹爪位置与旋转矩阵，
$\mathbf x_i^*,\mathbf R_i^*$ 为计划中的对应位姿。位置误差与旋转误差定义为
\[
 e_i^{\mathrm{pos}}=\|\mathbf x_i-\mathbf x_i^*\|_2,
 \qquad
 e_i^{\mathrm{rot}}=\arccos\!\left(\frac{\operatorname{tr}((\mathbf R_i^*)^{\mathsf T}\mathbf R_i)-1}{2}\right).
\]
数值计算将反余弦的输入截断到 $[-1,1]$。接近到位标签为
\[
 A_i=\mathbf{1}[e_i^{\mathrm{pos}}\leq\epsilon_p\ \land\ e_i^{\mathrm{rot}}\leq\epsilon_r],
 \qquad \epsilon_p=2\ \mathrm{mm},\quad\epsilon_r=0.5^\circ.
\]
设 $T_i$ 为排除初始化稳定阶段后的全部物理步，$f_i(t)$ 为该步记录的非夹指连杆接触合力与夹指非目标接触合力中的最大模长。
禁止接触约束的通过标签为
\[
 C_i=\mathbf{1}\!\left[\max_{t\in T_i}f_i(t)\leq F_{\max}\right],
 \qquad F_{\max}=0.5\ \mathrm N.
\]
令 $a_i$ 表示是否获得通过路径检查的已选计划。
有计划且执行记录完整时，任务成功标签为 $s_i=L_iA_iC_i$，对应原始记录中的 \texttt{execution\_success}；
无候选、无可执行计划或算法超时的任务取 $s_i=0$，物理夹持标签保持未测试。
基础设施中断单列为待补记录，补齐后再统计完整任务组。

设 $\mathcal I$ 为某方法与环境配置下预定的场景集合，$N=|\mathcal I|$，
$N_{\mathrm{plan}}=\sum_{i\in\mathcal I}a_i$，$N_{\mathrm{success}}=\sum_{i\in\mathcal I}s_i$。
全场景抓取成功率、路径通过率和已有计划的条件抓取成功率分别定义为
\[
 r_{\mathrm{grasp}}=\frac{N_{\mathrm{success}}}{N},
 \qquad r_{\mathrm{plan}}=\frac{N_{\mathrm{plan}}}{N},
 \qquad r_{\mathrm{grasp}\mid\mathrm{plan}}=\frac{N_{\mathrm{success}}}{N_{\mathrm{plan}}}.
\]
在全部计划均取得有效执行记录且 $N_{\mathrm{plan}}>0$ 时，三者满足
\[
 r_{\mathrm{grasp}}=r_{\mathrm{plan}}\,r_{\mathrm{grasp}\mid\mathrm{plan}}.
\]
主要结果报告 $r_{\mathrm{grasp}}$ 及其成功数与分母；$N_{\mathrm{plan}}=0$ 时条件成功率记为缺测。
候选数、空候选数、规划拒绝原因和执行失败阶段作为诊断记录。
放回精度也单独记录：最终平面位置与稳定初始位置的欧氏距离为 $e_i^{\mathrm{return}}$；
补充的放回通过标签要求 $L_i=1$、$e_i^{\mathrm{return}}\leq20$ mm，且放回保持阶段的各采样时刻均满足 $|h_i(t)|<3$ mm。
该标签不改变抓取成功率的定义。

几何指标先逐场景计算，再等权汇总。
对任一有定义的逐场景指标 $g_i$，令 $\mathcal I_g\subseteq\mathcal I$ 为其非缺测场景集合，均值为
\[
 \bar g=\frac{1}{|\mathcal I_g|}\sum_{i\in\mathcal I_g}g_i.
\]
缺测数与均值同时报告，点数另报中位数和范围。抓取成功率始终使用预定任务集合的完整分母。
两个配置在同一场景集合上的成功率差以百分点表示，定义为 $\Delta=100(r_{\mathrm{grasp}}^{(a)}-r_{\mathrm{grasp}}^{(b)})$。

\section{\qyb{RQ1: Quality of Stereo-Derived Point Clouds}}
\label{sec:rq1}

本节回答“\rqone”，确定进入抓取模型的几何信息具有怎样的误差、覆盖和目标纯度。

\subsection{\qyb{Motivation}}

感知阶段的异常会传递到候选生成和路径检查。
因此首先分析抓取模型实际收到的目标点云，明确候选生成前的几何输入质量。
这一分析为 RQ2 解释不同双目前端的下游表现提供场景级证据。

\subsection{\qyb{Analysis Method}}

对实验数据集全部 72 个场景的七种输出逐份核对点数、非有限坐标、单位、过滤前后变化和输入哈希。
三维误差采用预测点到目标真值网格的距离；近表面点占比统计给定距离内的预测点比例；
覆盖率以固定左右相机可见的真实目标表面采样点为分母。
阈值采用已保存的 1、2、5 mm 评价档位，正文用 2 mm 汇总，其余用于敏感性说明。
目标来源比例通过投回左图并读取离线实例身份计算，与表面距离分别报告。
全部目标有效像素的深度误差与仅保留目标点的误差分开，以反映过滤的选择作用。

\subsection{\qyb{Results and Analysis}}

\subsubsection{Point Count, Accuracy, and Coverage}

现有 504 份点云均满足接口校验，目标点数范围为 1,531--32,581，超过接口的 128 点下限。
\Cref{tab:rq1}说明满足点数要求后，各方法的近表面比例和可见覆盖仍存在差异。
结果图拟展示逐场景分布，并按同一工件并排显示左右图像、预测深度、目标点云和表面误差，
使用相同视角与色标，保留失败和异常场景。

% 数据：RQ1/data/development_pointclouds_20261012_v2/by_method.csv。
\begin{table}[H]
\centering\small
\caption{实验数据集的点云结果，每种方法 72 场景。点数为中位数，其余为逐场景指标的等权平均；覆盖率仅针对左右可见表面。}
\label{tab:rq1}
\begin{tabularx}{\linewidth}{Y r r r r}
\toprule
方法 & 目标点数 & 2 mm 内点（\%） & 2 mm 覆盖（\%） & 表面误差（mm） \\
\midrule
StereoBM & 7,367 & 31.6 & 67.8 & 30.30 \\
SGM（StereoSGBM） & 6,873 & 37.1 & 71.6 & 20.15 \\
PSMNet & 22,682 & 6.2 & 21.7 & 40.07 \\
GwcNet & 8,977 & 22.7 & 57.3 & 32.92 \\
RAFT-Stereo & 3,178 & 67.6 & 68.9 & 2.59 \\
IGEV-Stereo & 3,062.5 & 71.2 & 68.4 & 1.86 \\
FoundationStereo & 3,034 & 92.9 & 93.9 & 3.02 \\
\bottomrule
\end{tabularx}
\end{table}

\subsubsection{Target Extraction and Outlier Cases}

当前输出中，PSMNet 目标点来自真实目标像素的平均比例为 12.3\%，FoundationStereo 为 95.1\%，IGEV-Stereo 为 98.8\%。
固定空间区域与高度规则可能接收因深度误差移入区域的背景点，也可能接收邻件点。
因此分别分析目标像素深度误差和非目标来源，区分双目匹配与目标提取对输入质量的影响。
FoundationStereo 有 6/72 个场景的平均表面误差超过 5 mm，最大为 52.47 mm；
该尾部现象有助于解释它的近表面比例与覆盖较高，但平均距离仍高于 IGEV-Stereo。

\subsection{\qyb{Answer to RQ1}}

实验结果表明，各方法点云均满足接口点数要求，但目标几何的误差和覆盖存在差异。
FoundationStereo 在本批次的 2 mm 近表面点占比和覆盖率较高，IGEV-Stereo 的平均表面距离较小；
若干方法的主要输入问题同时涉及深度误差与非目标点混入。
正式回答将概括各维度的配对差异及异常场景。

\section{\qyb{RQ2: Effect of Stereo Perception on Grasping}}
\label{sec:rq2}

本节回答“\rqtwo”，比较七种双目感知方法在三个环境、三种场景条件下的抓取成功率。

\subsection{\qyb{Motivation}}

固定抓取模块可以观察替换双目前端后流程结果怎样变化。
这一比较包括共同目标提取和候选筛选的影响，能够检验几何指标较好的输入是否也带来更多可执行抓取。
失败分析以具体场景的点云、候选与路径记录为依据。

\subsection{\qyb{Analysis Method}}

读取七种前端在 72 个场景上的 504 个任务，固定 GraspGen 权重与种子；每个仿真环境内固定工具、初始位姿、路径检查和控制设置。
在 FANUC＋Isaac Sim、Franka＋Isaac Sim 和 Franka＋MuJoCo 中分别评价，每环境按单件、邻件、料框分层。
比较矩阵为七种感知方法、三种场景条件与三个环境；每格包含八件工件的三个位姿，共 24 个任务。
先报告有候选和通过路径检查的场景数，再关联各环境的全部物理结果；FANUC 有 156 条已选计划，两个 Franka 环境各有 160 条。
全场景抓取成功率以每方法、每环境的 72 场景为分母，分场景条件时以 24 为分母；条件成功率以该组已选计划数为分母。
默认与较低摩擦条件在各环境内保持原计划，材料敏感性作为本问的补充分析。

\subsection{\qyb{Results and Analysis}}

\subsubsection{From Point Clouds to Executable Plans}

候选生成、候选筛除与路径检查分别列出，拟按最先导致任务中止的阶段汇总互斥失败类别。
逐候选拒绝原因作为补充，因为同一任务可以有多个被拒绝候选。
路径数量差异结合目标点来源、夹爪开口、候选朝向、接触筛选与碰撞记录解释。
跨机械臂比较保留各自选中候选，避免把不同机器人可执行范围的差异归因于双目模型。

\subsubsection{Grasping Success in Three Environments}

\Cref{tab:rq2}按三种环境展示各双目方法的默认材料结果。
FoundationStereo 与 GraspGen 的成功数分别为 42/72、52/72 和 17/72，在当前三个环境中均高于其余双目前端组合。
方法间差值先在同一环境内比较，再检查差值方向和方法排序是否随环境变化。

% 数据：physics_validation_20261013/analysis/by_rq_method.csv 与 RQ2/results/platform_comparison_20261014/by_method.csv；nominal。
\begin{table}[H]
\centering\small
\caption{RQ2 默认材料下的抓取成功结果，固定 GraspGen。每种方法、每个环境均有 72 场景，包含三种场景条件各 24 个。每格给出成功数／总场景数及成功率，抓取成功需满足抬升保持、到位与禁止接触判定。}
\label{tab:rq2}
\begin{tabularx}{\linewidth}{Y r r r}
\toprule
双目方法 & FANUC＋Isaac Sim & Franka＋Isaac Sim & Franka＋MuJoCo \\
\midrule
StereoBM & 0/72（0.0\%） & 0/72（0.0\%） & 0/72（0.0\%） \\
SGM（StereoSGBM） & 0/72（0.0\%） & 0/72（0.0\%） & 0/72（0.0\%） \\
PSMNet & 0/72（0.0\%） & 0/72（0.0\%） & 0/72（0.0\%） \\
GwcNet & 2/72（2.8\%） & 3/72（4.2\%） & 0/72（0.0\%） \\
RAFT-Stereo & 17/72（23.6\%） & 16/72（22.2\%） & 7/72（9.7\%） \\
IGEV-Stereo & 18/72（25.0\%） & 20/72（27.8\%） & 8/72（11.1\%） \\
FoundationStereo & 42/72（58.3\%） & 52/72（72.2\%） & 17/72（23.6\%） \\
\bottomrule
\end{tabularx}
\end{table}

\subsubsection{Scene Conditions and Failure Localization}

\Cref{tab:rq2-scenes}汇总三种场景条件，各环境每种条件包含七种方法的 168 个任务。
正式分析按每个环境绘制七种方法在三种场景条件下的成功率，每个“方法--条件”单元以 24 为分母。
汇总表用于描述当前任务矩阵，场景难度比较保持方法、工件和位姿匹配。

% 数据：旧平台 per_candidate.csv 与 RQ2/data/platform_comparison_20261014/per_task.csv；nominal，graspgen。
\begin{table}[H]
\centering\small
\caption{RQ2 默认材料下按场景条件汇总的抓取成功计数。每格分母为七种方法各 24 个任务；同一场景被不同方法使用，统计推断保留这种配对关系。}
\label{tab:rq2-scenes}
\begin{tabularx}{\linewidth}{Y r r r}
\toprule
场景条件 & FANUC＋Isaac Sim & Franka＋Isaac Sim & Franka＋MuJoCo \\
\midrule
单件 & 20/168 & 25/168 & 8/168 \\
邻件 & 33/168 & 34/168 & 13/168 \\
料框 & 26/168 & 32/168 & 11/168 \\
\bottomrule
\end{tabularx}
\end{table}

按表中环境顺序，默认材料下成功数为 79/504、91/504、32/504，较低摩擦下为 49/504、52/504、22/504。
FoundationStereo 与 GraspGen 在 Franka＋Isaac Sim 的条件成功率为 52/52，全场景成功率为 52/72，两个分母分别报告。
对已到达抓取位姿但未抬升的场景，检查实际接触、闭合间隙和目标相对运动；
对抬升成功但执行未通过的场景，检查禁止接触的对象与发生阶段。
几何误差与成功标签按方法、工件和场景条件分层开展关联分析。

\subsection{\qyb{Answer to RQ2}}

在固定 GraspGen 和当前规划设置下，不同双目前端的抓取成功率随方法及环境变化。
FoundationStereo 组合在三个环境中取得较多成功场景，但绝对成功率分别为 58.3\%、72.2\% 和 23.6\%。
完整回答同时列出三种场景条件下的逐方法结果，并用点云、候选、路径和接触证据解释差异。

\section{\qyb{RQ3: Comparison of Grasp Generation Modules}}
\label{sec:rq3}

本节回答“\rqthree”，比较相同 FoundationStereo 输入下五种抓取模型在三个环境、三种场景条件下的成功率。

\subsection{\qyb{Motivation}}

即使点云完全相同，不同模型也会给出不同朝向、开口和分布的候选。
这些候选在统一夹爪和接近限制下的可用性，与原模型的内部评分是不同的评价对象。
本问关注现有模型作为可替换模块接入工业机械臂流程后的实际表现。

\subsection{\qyb{Analysis Method}}

在相同 72 个 FoundationStereo 点云上比较五个模型，共 360 个前端任务，并分别在三个仿真环境中评价。
比较矩阵为五种抓取模型、单件／邻件／料框三种场景条件和三个环境，每个单元有 24 个任务。
固定感知结果，每个环境内固定世界坐标、工具、路径检查和执行协议，核对各适配器的轴定义、宽度单位、工具偏移与候选排序。
记录模型原生要求及共同约束造成的筛除，区分空候选、运行异常和没有通过规划的候选。
FANUC 有 101 条已选计划，两个 Franka 环境各有 100 条，均已有物理结果。
GraspGen 在各环境中的任务直接复用 RQ2 的对应证据；全场景成功率的分母为每模型 72，分场景条件为每模型 24。

\subsection{\qyb{Results and Analysis}}

\subsubsection{Candidate Availability and Plan Acceptance}

先记录各模型候选及规划结果，再用\cref{tab:rq3}展示三个环境的全场景成功率。
6-DOF GraspNet 与 Contact-GraspNet 通过路径检查的计划较少，
需要先查明其候选如何被方向、开口、接触与运动约束筛除，再讨论模型的接入适用性。
补充图拟展示原始候选方向分布、逐规则筛除比例和代表场景。

\subsubsection{Grasping Success in Three Environments}

默认材料下，GraspGen 在三个环境分别成功 42/72、52/72 和 17/72；GPD 分别成功 18/72、19/72 和 2/72。
这些结果同时保留全场景分母和环境身份，用于分析固定感知后抓取模型的接入表现。

% 数据：physics_validation_20261013/analysis/by_rq_method.csv 与 RQ3/results/platform_comparison_20261014/by_method.csv；nominal。
\begin{table}[H]
\centering\small
\caption{RQ3 默认材料下的抓取成功结果，固定 FoundationStereo。每模型、每环境均有 72 场景，包含三种场景条件各 24 个。每格给出成功数／总场景数及成功率；GraspGen 与 RQ2 共享对应环境的任务和物理记录。}
\label{tab:rq3}
\begin{tabularx}{\linewidth}{Y r r r}
\toprule
抓取模型 & FANUC＋Isaac Sim & Franka＋Isaac Sim & Franka＋MuJoCo \\
\midrule
GPD & 18/72（25.0\%） & 19/72（26.4\%） & 2/72（2.8\%） \\
6-DOF GraspNet & 1/72（1.4\%） & 1/72（1.4\%） & 0/72（0.0\%） \\
Contact-GraspNet & 2/72（2.8\%） & 2/72（2.8\%） & 1/72（1.4\%） \\
GraspNet-Baseline & 1/72（1.4\%） & 1/72（1.4\%） & 1/72（1.4\%） \\
GraspGen & 42/72（58.3\%） & 52/72（72.2\%） & 17/72（23.6\%） \\
\bottomrule
\end{tabularx}
\end{table}

\subsubsection{Scene Conditions and Failure Localization}

\Cref{tab:rq3-scenes}列出三种场景条件的合并计数；逐模型分析按“模型 × 场景条件 × 环境”展示成功率，每个单元以 24 为分母。
先在同一环境内比较模型，再观察同一模型的结果随邻件、料框与环境配置怎样变化。

% 数据：旧平台 per_candidate.csv 与 RQ3/data/platform_comparison_20261014/per_task.csv；nominal，foundationstereo。
\begin{table}[H]
\centering\small
\caption{RQ3 默认材料下按场景条件汇总的抓取成功计数。每格包含五种模型各 24 个任务；推断与失败分析按模型、工件和对应位姿配对。}
\label{tab:rq3-scenes}
\begin{tabularx}{\linewidth}{Y r r r}
\toprule
场景条件 & FANUC＋Isaac Sim & Franka＋Isaac Sim & Franka＋MuJoCo \\
\midrule
单件 & 21/120 & 24/120 & 6/120 \\
邻件 & 29/120 & 32/120 & 8/120 \\
料框 & 14/120 & 19/120 & 7/120 \\
\bottomrule
\end{tabularx}
\end{table}

按表中环境顺序，默认材料下成功数为 64/360、75/360、21/360，较低摩擦下为 41/360、44/360、20/360。
失败分析分别核对空候选、规划筛除、接近到位、闭合接触、抬升保持及禁止碰撞。
两个 Franka 环境中的不一致结果可追溯到同一关节计划；涉及 FANUC 的比较还需考虑机器人结构、候选选择及驱动设置。
GraspDataGen 孤立夹爪抗扰动记录作为补充诊断，结合支承、动态接触与受力方向解释夹持结果。

\subsection{\qyb{Answer to RQ3}}

在固定 FoundationStereo 和当前工具、规划设置下，不同抓取模型在三个环境中的成功率存在差异。
GraspGen 的默认全场景成功率分别为 58.3\%、72.2\% 和 23.6\%，GPD 分别为 25.0\%、26.4\% 和 2.8\%。
其他模型在候选到已选计划阶段保留下来的任务较少，这一阶段的筛除与最终成功率一并分析。
正式结论同时报告三种场景条件下的逐模型结果，说明原始输出、适配、筛选和环境配置的作用，并用逐场景证据支持失败归因。

\section{\qyb{Discussion}}
\label{sec:discussion}

\subsection{\qyb{对工业界的启示}}

工业部署应根据各阶段证据定位抓取失败，再选择需要调整的模块。
目标点来源比例偏低时，优先验证预测目标分割和桌面去除；
候选在共同规则中大量被筛除时，核查夹爪坐标、开口语义和模型原生夹爪条件；
路径形式限制较强时，另设绕障规划与接近方向对照；
到位后滑落或碰撞时，结合接触记录检查工具几何、驱动与材料假设。
对每条失败任务记录最早阻断阶段，同时保留闭合、碰撞、宽度与超时等并存事件。
三个环境已执行的计划均达到末端到位要求，但夹持结果不同，后续排查重点应落在接触、工具及物理参数上。
上述改进在独立批次中评价，真实控制器接口与硬件通信另行验证。

方法选择应以工业工位的全场景抓取成功率及失败类型为依据。
同一配置分别报告单件、邻件和料框条件，每种条件均以 24 个场景为分母，观察遮挡、接近空间和壁面对结果的影响。
环境间分歧较大的方案应优先检查夹爪碰撞形状、摩擦、驱动与工件质量，并用真实工件进行小规模验证后确定部署配置。

\subsection{\qyb{对学术界的启示}}

抓取方法研究需要同时报告几何输入、候选、计划和物理抓取结果，明确阶段指标与最终成功率的关系。
全部场景的成功率保留无法生成计划的情况，已有计划的条件成功率用于分析执行阶段。
跨环境实验复用前端输入与候选，在运动学模型相同时进一步复用计划，公开逐任务配对记录及不一致案例，使算法变化与环境变化能够分别核查。
同一工件的重复位姿、共享方法组合和多个仿真结果按任务标识配对，并按工件聚合。

后续研究优先采用能够区分竞争解释的受控实验。
点云侧比较现有目标提取与冻结的预测分割，真值点云或真值掩码只作为单列的诊断上界；
抓取侧比较同一候选集合在原路径限制与新增绕障搜索下的差异；
执行侧在固定候选下检查时间步、摩擦和夹指驱动的敏感性。
后续可增加简单几何夹持启发式作为任务基线，在相同接口、预算和评价条件下比较。
由本次结果提出的改进应在新增独立工件上验证，保持原始配置与失败记录可复查。

\subsection{\qyb{Threats to Validity}}
\label{sec:threats}

\subsubsection{Internal Validity}

坐标、单位、夹爪偏移、模型适配及缓存对应错误会改变阶段结果。
采用带哈希的输入输出关联、关节与末端位姿核对、正负校准例及原始接触轨迹复算控制这些风险。
共同目标区域和朝下约束可能对不同方法产生不同影响，后续通过独立消融检查其作用。
物理参数和代码的修正保留版本，校准、实现诊断及主要实验分别计数。

\subsubsection{External Validity}

当前数据只有两类程序化工业形状，表面为哑光渲染，尚未覆盖真实金属反光、传感器噪声、透明材料和随机堆叠。
已归档预训练模型在本任务上保持冻结，其训练数据来源与合成资产的潜在相似性需要在模型清单中核对。
FANUC 与 Franka 模型、平行两指工具、统一质量和未由真机辨识的驱动参数限定了仿真结果的外推范围。
扩展工件、预测感知噪声和真实硬件测试是检验部署效果的后续步骤。

\subsubsection{Construct Validity}

点数、表面误差、候选评分、路径通过和物理任务完成测量不同对象，论文分别定义其分母与用途。
可见覆盖只针对左右视角，完整执行的成功条件同时包含抬升和接触门禁。
当前接触合力可能因相反方向的接触而抵消，后续应补充逐接触点检查及自碰覆盖审计。
GraspDataGen 抗扰动、单独物理抬升和本文规定的抓取成功使用不同判定条件，分析中分别报告。
三个环境同时涉及机器人结构、驱动、碰撞资产与接触求解差异，结果按完整环境配置解释。

\subsubsection{Conclusion Validity}

同一工件的多个位姿、不同种子和共享点云之间存在相关性，方法结果按场景配对、按工件及形状来源聚合。
当前工件覆盖两个形状族，类别间比较采用描述性统计。
后续扩展研究应增加独立几何来源、冻结随机种子与预算，报告配对效应和不确定性，并控制多重比较。
依据已观察结果调整的问题与参数按探索性分析报告，并记录修订时间与数据使用范围。

\section{\qyb{Conclusion}}
\label{sec:conclusion}

本文通过三个研究问题比较了双目点云质量、固定 GraspGen 时的感知方法表现，以及固定 FoundationStereo 时的抓取模型表现。
72 个合成工业场景的结果表明，点数充足的点云仍可能存在表面误差、覆盖不足和非目标点混入；
在给定工具与规划设置下，感知和抓取模块的选择均影响最终成功率。
FoundationStereo 与 GraspGen 组合在默认材料条件下，在 FANUC＋Isaac Sim、Franka＋Isaac Sim 和 Franka＋MuJoCo 中的成功率分别为
58.3\%、72.2\% 和 23.6\%，表明同一模块组合的执行表现随实验环境配置而变化。
因此，面向工业机械臂的方法选择应联合考察点云质量、路径可执行性和物理抓取结果，并保留逐任务的阶段记录以定位失败原因。

未来工作将扩展独立工业工件和场景，纳入真实金属表面、相机噪声及更复杂的遮挡条件，并在冻结配置后开展重复评价。
进一步搭建标定双目相机与工业机械臂组成的硬件环境，测量夹爪、摩擦和控制参数，
对照仿真与真实抓取的逐任务结果，检验方法差异能否在真实工位中保持。
针对已发现的失败类型，将分别评价目标点提取、候选适配和运动规划的改进，使各模块的调整具有独立的实验依据。

\begingroup
\small
\bibliographystyle{IEEEtran}
\bibliography{references}
\endgroup

\end{document}
